\chapter{1991:Richard R. Ernst}

\label{chap:chemistry1991}

The Nobel Prize in Chemistry for the year 1991 was awarded to Richard R. Ernst.

\textbf{Motivation:}

for his contributions to the development of the methodology of high resolution nuclear magnetic resonance (NMR) spectroscopy

\section{Richard R. Ernst}

\subsection*{Early Life and Education}

Richard R. Ernst was born in 1933 in Winterthur, Switzerland.

\subsection*{Career and Major Research}

Ernst earned his Ph.D. from the ETH Zurich in 1962, and he worked at Varian Associates in Palo Alto from 1962 to 1968. He returned to the ETH Zurich as a professor in 1968, where he spent the rest of his career.

\subsection*{Nobel Prize-winning Work}

The work for which Richard R. Ernst was awarded the Nobel Prize in Chemistry in 1991 was described by the Nobel Committee as: "for his contributions to the development of the methodology of high resolution nuclear magnetic resonance (NMR) spectroscopy".

He introduced Fourier transform NMR, which made spectra dramatically more sensitive, and he developed two-dimensional NMR, which separates overlapping signals into two frequency dimensions.

\subsection*{Significance and Impact}

These innovations transformed NMR into a general tool of chemistry and made possible the determination of structures of complex molecules in solution.

\subsection*{Later Career and Honors}

Ernst remained at the ETH Zurich, where he was appointed an emeritus professor. He received the 1991 Nobel Prize in Chemistry, and he died in 2021.

\subsection*{Legacy}

Fourier transform and two-dimensional NMR are the basis of modern NMR spectroscopy in chemistry, biology, and medicine.

\subsection*{Modern Developments}

Ernst's development of Fourier transform and two-dimensional NMR spectroscopy revolutionized molecular structure determination. Modern NMR, including multidimensional and solid-state methods, is now indispensable in chemistry, biology, and medicine, from the structure determination of proteins to magnetic resonance imaging in hospitals worldwide.

\subsection*{Selected Publications}

"Application of Fourier Transform Spectroscopy to Magnetic Resonance" (Review of Scientific Instruments, 1966)

\clearpage

\section*{Chapter Summary}

The 1991 Nobel Prize in Chemistry recognized methodological advances in high-resolution NMR spectroscopy. Ernst's Fourier transform and two-dimensional techniques increased the sensitivity and resolving power of NMR to the point that the structures of complex biomolecules could be determined in solution.
